% TOP_PROBLEM: {"problemId": "problem.beal-s-conjecture-tijdeman-zagier-conjecture", "canonicalTitle": "Beal's Conjecture (Tijdeman-Zagier Conjecture)", "releaseRank": 366, "primaryDomain": "number_theory_arithmetic_geometry", "primaryDomainLabel": "Number theory & arithmetic geometry", "displayStatus": "open", "releaseStatus": "open"}
% !TeX program = lualatex
% Definition and sources only; this file is not an LLM solution attempt.
\documentclass[11pt]{article}
\usepackage[margin=25mm]{geometry}
\usepackage{fontspec}
\setmainfont{Latin Modern Roman}
\usepackage{amsmath,amssymb,mathtools}
\usepackage{unicode-math}
\setmathfont{Latin Modern Math}
\usepackage{xurl,hyperref}
\hypersetup{colorlinks=true,urlcolor=blue}
\setcounter{secnumdepth}{0}
\setlength{\parindent}{0pt}
\setlength{\parskip}{5pt}
\setlength{\emergencystretch}{3em}
\begin{document}
\section*{\texorpdfstring{Beal's Conjecture (Tijdeman-Zagier Conjecture)}{Beal's Conjecture (Tijdeman-Zagier Conjecture)}}\label{366}
\subsection{Definitions and mathematical statement}
Beal's conjecture asserts that for positive integers $A,B,C$ and exponents $x,y,z>2$,
\[
 A^x+B^y=C^z\quad\Longrightarrow\quad\gcd(A,B,C)>1.
\]
Equivalently, some prime must divide all three bases. If a prime divides two bases in such an equation, it divides the third, so excluding a common prime also enforces pairwise coprimality. The exponents may be different and need not be prime. Positivity excludes zero and signed degeneracies. Unlike a finiteness statement for exceptional solutions, the assertion prohibits every coprime solution in this exponent range.
\par\smallskip\textit{Formulation and concept references:} \href{https://sites.math.unt.edu/~mauldin/beal.html}{[S1]}.

\subsection{Short English statement}
Whenever the sum of two positive perfect powers equals a third and all exponents exceed two, must their bases share a prime factor?
\subsection{Sources}
\begin{itemize}
\item[S1]R. D. Mauldin, The Beal Conjecture and Prize, University of North Texas; announced in Notices Amer. Math. Soc. 44 (1997), 1436-1437.
\url{https://sites.math.unt.edu/~mauldin/beal.html}.
\textit{Formulation source.}
\end{itemize}
\end{document}
